% Generated draft: requires mathematical and source review.
\documentclass[11pt]{article}
\usepackage[margin=25mm]{geometry}
\usepackage{fontspec}
\setmainfont{Latin Modern Roman}
\usepackage{amsmath,amssymb,mathtools,unicode-math}
\setmathfont{Latin Modern Math}
\usepackage{enumitem,needspace,xurl,hyperref}
\setlength{\parindent}{0pt}
\setlength{\parskip}{5pt}
\setlength{\emergencystretch}{3em}
\newcommand{\N}{\mathbb{N}}
\newcommand{\Z}{\mathbb{Z}}
\newcommand{\Q}{\mathbb{Q}}
\newcommand{\R}{\mathbb{R}}
\newcommand{\C}{\mathbb{C}}
\newcommand{\F}{\mathbb{F}}
\newcommand{\A}{\mathbb{A}}
\newcommand{\PP}{\mathbb{P}}
\newcommand{\E}{\mathbb{E}}
\newcommand{\eps}{\varepsilon}
\newcommand{\dd}{\,\mathrm{d}}
\begin{document}
\section*{UnsolvedMath ID: 30005194}
\textbf{Source problem code:} catalogue entry
\subsection{Definitions and mathematical statement}
A pair of \(C^*\)-algebras means an ordered pair \((A,B)\); any inclusion or additional structure required of this pair must be specified by the reconstruction theorem referenced below. For a groupoid \(G\) with unit space \(G^{(0)}\), source map \(s\), and range map \(r\), the isotropy group at \(x\in G^{(0)}\) is
\[
G_x^x=\{g\in G:s(g)=r(g)=x\}.
\]
The supplied passage refers to a reconstruction result called ``Theorem 1,'' but does not include its statement. Let \(\mathcal G_1\) denote the class of groupoids specified in that theorem, and let \(\mathcal P_1\) denote the collection of \(C^*\)-algebra pairs appearing there. These are source-dependent placeholders, not independently specified classes.

The source asks for the following three related developments:
\begin{enumerate}
\item Characterize \(\mathcal P_1\) by intrinsic \(C^*\)-algebraic properties. The requested characterization has the schematic form
\[
(A,B)\in\mathcal P_1
\quad\Longleftrightarrow\quad
(A,B)\text{ satisfies the proposed intrinsic conditions},
\]
with the admissible pairs and their additional structure determined by Theorem 1.
\item Obtain a version of Theorem 1 for twisted \(C^*\)-algebras over the same class \(\mathcal G_1\) of groupoids.
\item Determine whether such a twisted version can be used to weaken the hypotheses on the isotropy groups in Theorem 1.
\end{enumerate}
The supplied summary describes the setting as quasi-Cartan groupoid reconstruction. However, neither the meaning of quasi-Cartan in this setting, the reconstruction construction and conclusion, the class \(\mathcal G_1\), the permitted twists and associated \(C^*\)-algebra completion, nor the existing isotropy hypotheses are supplied. Consequently, this is a source-dependent statement of the research targets; a fully specified quantified formulation requires the missing theorem and its definitions.
\subsection{Short English statement}
Characterize the \(C^*\)-algebra pairs in the referenced quasi-Cartan reconstruction theorem, extend that theorem to twisted groupoid \(C^*\)-algebras over its original groupoid class, and determine whether the extension permits weaker isotropy assumptions. Theorem 1 is absent from the supplied record, so its exact hypotheses and conclusions remain unresolved.
\subsection{Sources}
\begin{itemize}
\item Local UnsolvedMath registry, permanent ID \(30005194\), source problem code \texttt{catalogue entry}, title ``Twisted Groupoid Reconstruction of \(C^*\)-Algebra Pairs.'' The supplied record belongs to the Oberwolfach Reports Open Problems collection. No individual proposer is named.
\item Supplied report citation, with mathematical typography normalized: \(C^*\)-Algebras (2023), \emph{Oberwolfach Reports}, DOI 10.4171/owr/2022/36, pp.\ 2059--2127. Original report: \url{https://doi.org/10.4171/owr/2022/36}. The extracted passage refers to ``Theorem 1'' without reproducing it. The registry separately records proposed year 2022.
\item Supplied literature source: ``Cartan subalgebras for non-principal twisted groupoid C-star algebras,'' \url{https://arxiv.org/abs/2001.08270}.
\item Literature assessment dated exactly \texttt{2026-08-22}: the supplied curation labels the combined entry open, reports that the report explicitly raises these three questions, and states that later work supplies sufficient cases without a complete answer to the combined entry. This is supplied dated curation, not independently verified here; time-sensitive status and cited claims require independent verification.
\item The registry labels the statement \texttt{corrected\_verified} and states that it was reconstructed and checked against the official report. These are supplied provenance claims, not verification performed for this draft.
\end{itemize}
\end{document}
